% Generated draft: requires mathematical and source review.
\documentclass[11pt]{article}
\usepackage[margin=25mm]{geometry}
\usepackage{fontspec}
\setmainfont{Latin Modern Roman}
\usepackage{amsmath,amssymb,mathtools,unicode-math}
\setmathfont{Latin Modern Math}
\usepackage{enumitem,needspace,xurl,hyperref}
\setlength{\parindent}{0pt}
\setlength{\parskip}{5pt}
\setlength{\emergencystretch}{3em}
\newcommand{\N}{\mathbb{N}}
\newcommand{\Z}{\mathbb{Z}}
\newcommand{\Q}{\mathbb{Q}}
\newcommand{\R}{\mathbb{R}}
\newcommand{\C}{\mathbb{C}}
\newcommand{\F}{\mathbb{F}}
\newcommand{\A}{\mathbb{A}}
\newcommand{\PP}{\mathbb{P}}
\newcommand{\E}{\mathbb{E}}
\newcommand{\eps}{\varepsilon}
\newcommand{\dd}{\,\mathrm{d}}
\begin{document}
\section*{UnsolvedMath ID: 10400196}
\textbf{Source problem code:} catalogue entry
\subsection{Definitions and mathematical statement}
For an abelian group \(G\), the source calls a function \(q\) quadratic when its polarization
\[
 b_q(x,y)=q(x+y)-q(x)-q(y)
\]
is bilinear. It is homogeneous if \(q(nx)=n^2q(x)\) for every \(n\in\Z\) and \(x\in G\). The excerpt does not give the value group or identify \(G\).

It reports a canonical assignment \(\sigma\mapsto q_\sigma\) from \(\operatorname{spin}^c\) structures on the intended 3-manifolds to quadratic functions, with \(q_\sigma\) homogeneous if and only if \(\sigma\) comes from a spin structure. Not every \(\operatorname{spin}^c\) structure is induced by a spin structure; the question is not restricted to the homogeneous case.

F. Deloup asks whether \(\arg\gamma(q_\sigma)\) admits a lift to a mod-16 invariant. The source states that such a lift would give a finite-type invariant of degree one in the \(\operatorname{spin}^c\) Goussarov--Habiro theory. This is the requested lift and its reported consequence, not a claimed existing invariant.

The function \(\gamma\), its Gauss-sum normalization and nonvanishing conditions, the scaling of \(\arg\), and the codomain and reduction map for ``mod 16'' are missing. An ordinary argument of a nonzero complex number is valued in \(\R/(2\pi\Z)\); identifying it with a mod-8 class, and then choosing a reduction from a mod-16 target, requires extra conventions. In particular a \(\Z/16\Z\) target is not silently assumed when rational-valued refinements may be intended. The underlying quadratic refinement, manifold category and degree-one surgery conventions also need the surrounding source.
\subsection{Short English statement}
Can the Gauss-phase quantity attached to a \(\operatorname{spin}^c\) quadratic function be lifted to a mod-16 invariant, yielding the stated degree-one finite-type invariant? The phase normalization and lift map are absent, and nonhomogeneous \(\operatorname{spin}^c\) cases must be retained.
\subsection{Sources}
\begin{itemize}
\item Local UnsolvedMath registry, permanent ID 10400196, source code catalogue entry; statement attribution F. Deloup; proposer field and proposal year absent; update 2026-08-06.
\item Ohtsuki, \emph{Problems on invariants of knots and 3-manifolds} (2002), Question 10.21, PDF page 154, including the supplied quadratic-function footnote 37: \url{https://msp.org/gtm/2002/04/gtm-2002-04-024s.pdf}. Registry access date 2026-07-29; PDF font-structure extraction. Provenance distinguishes 231 numbered bold headings. Present primary-literature status and redistribution terms require review.
\item Adjacent local ID 10400195 concerns closed \(\operatorname{spin}^c\) 3-manifolds and \(Y_d^c\)-equivalence, but does not restore the missing phase or quadratic data. The supplied summary gives no exact resolving reference.
\end{itemize}
\end{document}
